\documentclass[a4paper,11pt]{article}
\usepackage[utf8]{inputenc}
\usepackage{geometry}
\geometry{left=2.5cm,right=2.5cm,top=2.5cm,bottom=2.5cm}
\usepackage{amsmath,amssymb}
\usepackage{booktabs}
\usepackage{graphicx}
\usepackage{float}
\usepackage{caption}
\usepackage{subcaption}
\usepackage{hyperref}
\usepackage{listings}
\usepackage{xcolor}
\usepackage{enumitem}
\usepackage{multirow}

\definecolor{codegreen}{rgb}{0,0.6,0}
\definecolor{codegray}{rgb}{0.5,0.5,0.5}
\definecolor{codepurple}{rgb}{0.58,0,0.82}
\definecolor{backcolour}{rgb}{0.95,0.95,0.92}

\lstdefinestyle{mystyle}{
    backgroundcolor=\color{backcolour},   
    commentstyle=\color{codegreen},
    keywordstyle=\color{magenta},
    numberstyle=\tiny\color{codegray},
    stringstyle=\color{codepurple},
    basicstyle=\ttfamily\footnotesize,
    breakatwhitespace=false,         
    breaklines=true,                 
    captionpos=b,                    
    keepspaces=true,                 
    numbers=left,                    
    numbersep=5pt,                  
    showspaces=false,                
    showstringspaces=false,
    showtabs=false,                  
    tabsize=2
}

\lstset{style=mystyle}

\title{Diabetes Prediction Using Machine Learning \\
\large A Comprehensive XGBoost-Based Approach with Cross-Validation}
\author{Your Name \\
\texttt{student.id@university.edu}}
\date{\today}

\begin{document}

\maketitle

\begin{abstract}
This project presents a machine learning approach for predicting diabetes using the Kaggle Diabetes Prediction Challenge dataset. The dataset contains 700,000 training samples and 300,000 test samples with 26 clinical and demographic features. We develop a systematic modeling pipeline that progresses from baseline models to an optimized XGBoost model with proper cross-validation. Key contributions include: (1) comprehensive exploratory data analysis revealing class imbalance and feature correlations; (2) progressive model development from simple baselines to advanced gradient boosting; (3) critical methodological improvements including 5-fold stratified cross-validation, fold-wise categorical encoding, and out-of-fold evaluation; (4) feature engineering based on clinical domain knowledge; and (5) model interpretation using feature importance analysis. The final XGBoost model achieves a mean cross-validation AUC of 0.7256 with low variance (0.0008), demonstrating stable generalization performance. The complete pipeline is implemented in a modular Python package with full reproducibility.
\end{abstract}

\section{Introduction}

\subsection{Background}
Diabetes is a chronic metabolic disorder affecting millions worldwide. Early detection and risk assessment are crucial for timely intervention. Machine learning models can assist healthcare professionals by identifying high-risk individuals based on clinical measurements and demographic information.

\subsection{Problem Statement}
Given a dataset of clinical measurements and demographic information, the task is to predict whether an individual has been diagnosed with diabetes. The evaluation metric is ROC-AUC, which measures the model's ability to distinguish between positive and negative cases across all classification thresholds.

\subsection{Dataset Description}
The dataset is synthetically generated from a deep learning model trained on the Diabetes Health Indicators Dataset. It consists of:
\begin{itemize}
    \item Training set: 700,000 samples, 26 features
    \item Test set: 300,000 samples, 25 features (target excluded)
    \item Target variable: \texttt{diagnosed\_diabetes} (binary: 0 = no diabetes, 1 = diabetes)
\end{itemize}

\section{Exploratory Data Analysis}

\subsection{Target Distribution}
The target variable shows imbalance: 62.3\% of samples are diabetic (class 1) and 37.7\% are non-diabetic (class 0). This imbalance ratio of 1.65 indicates that class 1 is the majority.

\subsection{Feature Correlations}
Correlation analysis reveals that \texttt{family\_history\_diabetes} has the strongest positive correlation with the target (0.211), followed by \texttt{age} (0.161) and \texttt{systolic\_bp} (0.107). Physical activity shows a negative correlation (-0.170), consistent with medical knowledge.

\subsection{Categorical Features}
Six categorical features are present: gender, ethnicity, education level, income level, smoking status, and employment status. No missing values are present in the dataset.

\subsection{Key Observations}
\begin{itemize}
    \item No missing values, simplifying preprocessing
    \item Class imbalance requires careful evaluation strategy
    \item Strong predictive signals from family history and age
    \item Outliers exist in physical activity and sleep hours features
\end{itemize}

\section{Methodology}

\subsection{Baseline Models}

\subsubsection{Dummy Classifier}
A dummy classifier predicting the majority class (diabetic) establishes a lower bound for performance. This achieves an AUC of 0.5000, representing random classification.

\subsubsection{Decision Tree}
A decision tree with max depth 5 serves as an interpretable baseline. It achieves an OOF AUC of 0.6825, with feature importance dominated by \texttt{family\_history\_diabetes} (44.8\%), \texttt{physical\_activity} (29.9\%), and \texttt{age} (21.8\%).

\subsection{XGBoost Model Development}

\subsubsection{Initial Model}
The initial XGBoost model with default parameters (100 trees, learning rate 0.1) achieves AUC 0.7192, significantly outperforming the decision tree baseline.

\subsubsection{Model Capacity Tuning}
Increasing tree count to 300 improves AUC to 0.7255, while further increasing to 500 yields 0.7264. Learning curves show convergence beyond 500 trees, indicating diminishing returns.

\subsubsection{Feature Engineering}
Based on clinical domain knowledge, five engineered features are added:
\begin{itemize}
    \item \texttt{age\_family\_history}: interaction between age and family history
    \item \texttt{age\_bmi}: interaction between age and body mass index
    \item \texttt{cardio\_risk\_score}: composite cardiovascular risk score
    \item \texttt{cholesterol\_ratio}: total cholesterol divided by HDL
    \item \texttt{non\_hdl\_cholesterol}: total minus HDL cholesterol
\end{itemize}

\subsection{Critical Methodological Improvements}

\subsubsection{From Fixed Holdout to Cross-Validation}
Early experiments used a single fixed validation split (random\_state=42). This approach risks overfitting to that specific split. The final pipeline employs 5-fold stratified cross-validation, ensuring each sample is validated exactly once.

\subsubsection{Fold-Wise Categorical Encoding}
To prevent data leakage, categorical encoding is performed within each fold:
\begin{lstlisting}[language=Python]
def _encode_categorical_fold(X_train_raw, X_valid_raw, categorical_cols):
    for col in categorical_cols:
        le = LabelEncoder()
        le.fit(X_train_raw[col])      # Fit only on fold-train
        X_train_encoded[col] = le.transform(X_train_raw[col])
        X_valid_encoded[col] = le.transform(X_valid_raw[col])
\end{lstlisting}

\subsubsection{Out-of-Fold Evaluation}
OOF (Out-of-Fold) predictions collect validation predictions from all folds, providing an unbiased estimate of generalization performance.

\subsubsection{Early Stopping}
Training stops when validation AUC fails to improve for 100 consecutive rounds, automatically determining optimal tree count.

\subsection{Final Model Architecture}
The final XGBoost model uses optimized hyperparameters:
\begin{itemize}
    \item \texttt{n\_estimators}: 2000 (early stopping active)
    \item \texttt{learning\_rate}: 0.03
    \item \texttt{max\_depth}: 4
    \item \texttt{subsample}: 0.8
    \item \texttt{colsample\_bytree}: 0.8
    \item \texttt{gamma}: 1 (split regularization)
    \item \texttt{reg\_alpha}: 0.1, \texttt{reg\_lambda}: 1 (L1/L2 regularization)
\end{itemize}

\section{Experimental Results}

\subsection{Cross-Validation Performance}

Table \ref{tab:cv_results} shows the 5-fold cross-validation results.

\begin{table}[H]
\centering
\caption{5-Fold Cross-Validation Results}
\label{tab:cv_results}
\begin{tabular}{lcc}
\toprule
\textbf{Fold} & \textbf{AUC} & \textbf{Best Iterations} \\
\midrule
Fold 1 & 0.72596 & 1993 \\
Fold 2 & 0.72428 & 1997 \\
Fold 3 & 0.72515 & 1999 \\
Fold 4 & 0.72675 & 1998 \\
Fold 5 & 0.72591 & 1999 \\
\midrule
\textbf{Mean} & \textbf{0.72561} & \textbf{1997} \\
\textbf{Std} & \textbf{0.00084} & - \\
\textbf{OOF AUC} & \textbf{0.72560} & - \\
\bottomrule
\end{tabular}
\end{table}

The low standard deviation (0.00084) indicates consistent performance across folds, demonstrating strong generalization capability.

\subsection{Model Comparison}

Table \ref{tab:comparison} compares all models using the same CV protocol.

\begin{table}[H]
\centering
\caption{Model Performance Comparison}
\label{tab:comparison}
\begin{tabular}{lccc}
\toprule
\textbf{Model} & \textbf{Mean CV AUC} & \textbf{Std} & \textbf{OOF AUC} \\
\midrule
Dummy Classifier & 0.5000 & 0.0000 & 0.5000 \\
Decision Tree & 0.6825 & 0.0012 & 0.6824 \\
XGBoost (Final) & \textbf{0.7256} & \textbf{0.0008} & \textbf{0.7256} \\
\bottomrule
\end{tabular}
\end{table}

XGBoost achieves a 6.3\% absolute improvement over the decision tree baseline.

\subsection{Feature Importance Analysis}

Figure \ref{fig:feature_importance} displays the top 15 feature importances from the final model.

\begin{figure}[H]
\centering
\includegraphics[width=0.8\textwidth]{figures/02_feature_importance.png}
\caption{Top 15 Feature Importance}
\label{fig:feature_importance}
\end{figure}

The top 5 features and their importance scores are:
\begin{enumerate}
    \item \texttt{family\_history\_diabetes}: 42.7\%
    \item \texttt{age\_family\_history}: 29.4\%
    \item \texttt{age\_bmi}: 6.9\%
    \item \texttt{cardio\_risk\_score}: 4.6\%
    \item \texttt{physical\_activity\_minutes\_per\_week}: 3.5\%
\end{enumerate}

\textbf{Key insights:}
\begin{itemize}
    \item Family history is the strongest single predictor
    \item The interaction feature \texttt{age\_family\_history} contributes significantly, indicating age effects are amplified when family history exists
    \item Engineered features (\texttt{cardio\_risk\_score}) rank in top 5, validating domain knowledge integration
    \item Physical activity shows expected negative correlation with diabetes risk
\end{itemize}

\subsection{Model Diagnostics}

\subsubsection{ROC Curve}
Figure \ref{fig:roc} shows the ROC curve from OOF predictions, with AUC = 0.7256, significantly better than random classification.

\begin{figure}[H]
\centering
\includegraphics[width=0.6\textwidth]{figures/04_roc_curve.png}
\caption{ROC Curve (Out-of-Fold Predictions)}
\label{fig:roc}
\end{figure}

\subsubsection{Prediction Distribution}
Figure \ref{fig:distribution} shows the prediction distributions for diabetic and non-diabetic cases, demonstrating good class separation.

\begin{figure}[H]
\centering
\includegraphics[width=0.8\textwidth]{figures/05_prediction_distribution.png}
\caption{Prediction Distribution by Class}
\label{fig:distribution}
\end{figure}

\subsection{Kaggle Submission Results}

The final model was submitted to Kaggle, achieving:
\begin{itemize}
    \item Public Score: 0.69561
    \item Private Score: 0.69314
\end{itemize}

The slight drop from CV AUC (0.7256) to test score (0.6956) is expected due to the synthetic nature of the data and potential distribution shifts between training and test sets.

\section{Discussion}

\subsection{Key Findings}

\begin{enumerate}
    \item \textbf{Cross-validation is essential}: The transition from fixed holdout to 5-fold CV revealed consistent performance, eliminating overfitting concerns.
    
    \item \textbf{Feature engineering adds value}: Domain-informed features like \texttt{cardio\_risk\_score} entered the top 5 importance rankings, validating the approach.
    
    \item \textbf{Interaction effects matter}: The \texttt{age\_family\_history} interaction contributes nearly as much as the original family history feature, indicating non-linear relationships.
    
    \item \textbf{Model stability}: Low CV variance (0.0008) confirms the model's generalization capability.
\end{enumerate}

\subsection{Lessons Learned}

\subsubsection{What Worked}
\begin{itemize}
    \item Stratified cross-validation for reliable evaluation
    \item Fold-wise categorical encoding preventing data leakage
    \item Early stopping for automatic model complexity control
    \item Domain-informed feature engineering
\end{itemize}

\subsubsection{What Could Be Improved}
\begin{itemize}
    \item Further feature engineering (e.g., polynomial interactions)
    \item Hyperparameter search over larger space
    \item Ensemble methods combining XGBoost with LightGBM
    \item SHAP analysis for enhanced interpretability
\end{itemize}

\section{Conclusion}

This project developed a systematic machine learning pipeline for diabetes prediction, progressing from simple baselines to an optimized XGBoost model with proper cross-validation. The final model achieves a mean CV AUC of 0.7256 with low variance (0.0008), demonstrating stable generalization. Key methodological contributions include fold-wise categorical encoding, out-of-fold evaluation, and early stopping for automatic model selection. The complete pipeline is implemented in a modular Python package, ensuring reproducibility.

The project highlights the importance of rigorous evaluation protocols in machine learning and demonstrates how domain knowledge can be effectively integrated through feature engineering. While the Kaggle test score (0.6956) is slightly lower than CV performance due to data distribution shifts, the model's consistent cross-validation results and stable performance across folds provide confidence in its generalization capability.

\section*{Code Availability}

The complete project code is organized as follows:

\begin{lstlisting}[language=bash]
machinelearning_project/
├── data/                    # Dataset files
├── notebooks/               # Jupyter notebooks (EDA, experiments)
├── src/                     # Python source code
│   ├── config.py           # Configuration and paths
│   ├── features.py         # Feature engineering
│   ├── preprocessing.py    # Data preprocessing
│   ├── models.py           # Model training with CV
│   ├── train_xgb_cv.py     # Main training script
│   └── evaluate.py         # Evaluation utilities
├── models/                  # Saved models
├── submissions/             # Kaggle submission files
├── logs/                    # Experiment logs
└── figures/                 # Generated plots
\end{lstlisting}

\section*{References}

\begin{enumerate}
    \item Chen, T., \& Guestrin, C. (2016). XGBoost: A Scalable Tree Boosting System. In Proceedings of the 22nd ACM SIGKDD International Conference on Knowledge Discovery and Data Mining.
    
    \item Kaggle. (2024). Diabetes Prediction Challenge. \url{https://www.kaggle.com/competitions/diabetes-prediction-challenge}
    
    \item American Diabetes Association. (2023). Standards of Medical Care in Diabetes.
    
    \item Kohavi, R. (1995). A Study of Cross-Validation and Bootstrap for Accuracy Estimation and Model Selection. IJCAI.
    
    \item Lundberg, S. M., \& Lee, S. I. (2017). A Unified Approach to Interpreting Model Predictions. NeurIPS.
\end{enumerate}

\end{document}
